
\section{Conclusion}
\label{section: conclusion}
In this paper, we introduce LEGO-GraphRAG, a unified framework for the modular analysis and design of GraphRAG instances. By dividing the GraphRAG retrieval process into distinct modules and identifying corresponding design solutions, LEGO-GraphRAG provides a viable approach for building advanced GraphRAG systems.
Building upon the LEGO-GraphRAG framework, we conduct extensive empirical studies on large-scale real-world graphs and diverse GraphRAG query sets, leading to key findings:
{\bf 1)} GraphRAG must balance quality, efficiency, and cost across three aspects: modules, method types, and semantic models.
{\bf 2)} Extracting query-relevant subgraphs is the primary efficiency bottleneck for GraphRAG on large-scale graphs and remains underexplored.
{\bf 3)} Graph structural information enables efficient solutions for GraphRAG, while semantic information is crucial for improving the quality of complex queries. The optimal solution should integrate structural information to identify query-relevant subgraphs and leverage semantic information to retrieve reasoning paths.



